\documentclass[11pt]{article}
\usepackage{amsmath,amssymb,geometry}
\geometry{margin=22mm}
\emergencystretch=4em
\begin{document}
\section*{P1: Sample-average mixed-binary linear model for tandem clinic appointments}
\paragraph{Source and scope} Original paper, PDF pages 1, 6, 7, 8, 9, 10, 11, and 12, equations (4)--(9), Remark 1, and Proposition 2, equations (11)--(25).
\subsection*{Sets, parameters, and scope}
Each patient $k=1,\ldots,K$ visits every single-server station $n=1,\ldots,N$ in the same fixed order, and all patients attend. $\Omega$ is a given finite, equally weighted sample of service-duration scenarios. $\delta_k^n(\omega)\in\mathbb Z_+$ is the number of time slots, with service times allowed to be correlated across patients and stations. $T^n\in\mathbb Z_+$ is the regular working time at the station. The cost inequalities $0\leq C_{L,k}\leq C_{W,k}$, $0\leq C_R^n\leq C_I^n$, and $C_O^n\geq0$ correspond to leisure and regular waiting, usable and regular idle time, and overtime, respectively. A valid big-M value is
\[
M=\max\{1,\max_nT^n+\max_{\omega\in\Omega}\sum_{k=1}^K\sum_{n=1}^N\delta_k^n(\omega)\}.
\]
This file uses the sample-average mixed-binary linear formulation in Proposition 2. Appointment times $a_k^n$ are continuous. Remark 1 permits relaxing time-slot integrality, and Proposition 2 uses a continuous domain. The objective is average cost over the given sample of service durations.
\subsection*{Variables and complete model}
$a$ is the appointment schedule chosen in advance. In each scenario, $r$ is the patient-ready time, $v$ is the provider-ready time, $s$ is the actual service start time, and $f$ is the provider's finish time. The binary variables $y,z$ select the branches of the maximum relationships. $\Gamma_W,\Gamma_L,\Gamma_I,\Gamma_R,\Gamma_O$ are the total durations of regular waiting, leisure waiting, regular idle time, usable idle time, and overtime, respectively. Boundary variables $y,z$ that do not apply may be omitted.
\begin{equation}
\min
\frac{1}{|\Omega|} \sum_{\omega \in \Omega}
\Big[
\sum_{k=1}^{K}
\big(
C_{W,k}\Gamma_{W,k}(\omega) + C_{L,k}\Gamma_{L,k}(\omega)
\big)
+ \sum_{n=1}^{N}
\big(
C_I^n\Gamma_I^n(\omega) + C_R^n\Gamma_R^n(\omega) + C_O^n\Gamma_O^n(\omega)
\big)
\Big]
\tag{MILP}
\end{equation}

\begin{equation}
\Gamma_{W,k}(\omega) =
\sum_{n=1}^{N} \left( s_k^n(\omega) - r_k^n(\omega) \right),
\qquad \forall k,\ \omega
\tag{10a}
\end{equation}

\begin{equation}
\Gamma_{L,k}(\omega) =
\sum_{n=2}^{N} \left( r_k^n(\omega) - s_k^{n-1}(\omega) - \delta_k^{n-1}(\omega) \right),
\qquad \forall k,\ \omega
\tag{10b}
\end{equation}

\begin{equation}
\Gamma_I^n(\omega) =
\sum_{k=1}^{K} \left( s_k^n(\omega) - v_k^n(\omega) \right),
\qquad \forall n,\ \omega
\tag{10c}
\end{equation}

\begin{equation}
\Gamma_R^n(\omega) =
\sum_{k=2}^{K} \left( v_k^n(\omega) - s_{k-1}^n(\omega) - \delta_{k-1}^n(\omega) \right),
\qquad \forall n,\ \omega
\tag{10d}
\end{equation}

\begin{equation}
\Gamma_O^n(\omega) = f^n(\omega) - T^n,
\qquad \forall n,\ \omega
\tag{10e}
\end{equation}

\begin{equation}
s_k^n(\omega) \ge r_k^n(\omega),
\qquad \forall k,\ n,\ \omega
\tag{11}
\end{equation}

\begin{equation}
s_k^n(\omega) \ge v_k^n(\omega),
\qquad \forall k,\ n,\ \omega
\tag{12}
\end{equation}

\begin{equation}
r_k^1(\omega) = a_k^1,
\qquad \forall k,\ \omega
\tag{13a}
\end{equation}

\begin{equation}
r_k^n(\omega) \ge a_k^n,
\qquad \forall k,\ n \ge 2,\ \omega
\tag{13b}
\end{equation}

\begin{equation}
r_k^n(\omega) \ge s_k^{n-1}(\omega) + \delta_k^{n-1}(\omega),
\qquad \forall k,\ n \ge 2,\ \omega
\tag{14}
\end{equation}

\begin{equation}
r_k^n(\omega) \le a_k^n + M(1-z_k^n(\omega)),
\qquad \forall k,\ n \ge 2,\ \omega
\tag{15}
\end{equation}

\begin{equation}
r_k^n(\omega) \le s_k^{n-1}(\omega) + \delta_k^{n-1}(\omega) + M z_k^n(\omega),
\qquad \forall k,\ n \ge 2,\ \omega
\tag{16}
\end{equation}

\begin{equation}
v_1^n(\omega) = a_1^n,
\qquad \forall n,\ \omega
\tag{17}
\end{equation}

\begin{equation}
v_k^n(\omega) \ge a_k^n,
\qquad \forall k \ge 2,\ n,\ \omega
\tag{18}
\end{equation}

\begin{equation}
v_k^n(\omega) \ge s_{k-1}^n(\omega) + \delta_{k-1}^n(\omega),
\qquad \forall k \ge 2,\ n,\ \omega
\tag{19}
\end{equation}

\begin{equation}
v_k^n(\omega) \le a_k^n + M(1-y_k^n(\omega)),
\qquad \forall k \ge 2,\ n,\ \omega
\tag{20}
\end{equation}

\begin{equation}
v_k^n(\omega) \le s_{k-1}^n(\omega) + \delta_{k-1}^n(\omega) + M y_k^n(\omega),
\qquad \forall k \ge 2,\ n,\ \omega
\tag{21}
\end{equation}

\begin{equation}
f^n(\omega) \ge T^n,
\qquad \forall n,\ \omega
\tag{22}
\end{equation}

\begin{equation}
f^n(\omega) \ge s_K^n(\omega) + \delta_K^n(\omega),
\qquad \forall n,\ \omega
\tag{23}
\end{equation}

\begin{equation}
a_k^n \ge a_{k-1}^n,
\qquad \forall k \ge 2,\ n
\tag{24}
\end{equation}

\begin{equation}
a_k^n \ge a_k^{n-1},
\qquad \forall k,\ n \ge 2
\tag{25}
\end{equation}

\begin{equation}
a_1^1 = 0,\quad
0 \le a_k^n \le T^n,\quad
s_k^n(\omega),\ r_k^n(\omega),\ v_k^n(\omega),\ f^n(\omega) \ge 0,
\qquad \forall k,\ n,\ \omega
\tag{26a}
\end{equation}

\begin{equation}
\Gamma_{W,k}(\omega),\ \Gamma_{L,k}(\omega),\ \Gamma_I^n(\omega),\ \Gamma_R^n(\omega),\ \Gamma_O^n(\omega) \ge 0,
\qquad \forall k,\ n,\ \omega
\tag{26b}
\end{equation}

\begin{equation}
z_k^n(\omega),\ y_k^n(\omega) \in \{0,1\},
\qquad \forall k,\ n,\ \omega
\tag{26c}
\end{equation}


\end{document}
